\chapter{INTRODUCCIÓN}

En las últimas dos décadas, la informática ha dejado de ser un campo reservado exclusivamente a especialistas para convertirse en una competencia transversal con impacto en la vida cotidiana, la educación y el trabajo. Comprender cómo se representan los problemas, cómo se estructuran las soluciones y cómo se diseñan procesos ordenados para resolverlos constituye hoy una necesidad formativa cada vez más relevante, incluso para personas que no se dedicarán profesionalmente al desarrollo de software. En este contexto, el pensamiento computacional ha adquirido un papel central dentro de las discusiones educativas contemporáneas, al proponer una forma de razonamiento basada en la descomposición de problemas, el reconocimiento de patrones, la abstracción y el diseño de procedimientos paso a paso.

Sin embargo, la incorporación del pensamiento computacional en contextos escolares plantea un desafío importante: no basta con reconocer su valor teórico, sino que es necesario traducirlo en experiencias educativas accesibles, motivadoras y adecuadas a la edad de los estudiantes. En muchos casos, la enseñanza de conceptos vinculados a la informática se confunde con la enseñanza de programación o con el uso instrumental de herramientas digitales, lo que limita su alcance pedagógico. Frente a esta situación, han surgido iniciativas internacionales que buscan acercar los fundamentos de la informática a niños y jóvenes sin exigir conocimientos técnicos previos. Entre ellas, una de las más consolidadas es el Desafío Bebras.

El Desafío Bebras es una competencia internacional orientada a promover el pensamiento computacional mediante tareas breves que plantean problemas de lógica, representación, secuencias, estructuras, algoritmos y análisis de situaciones. Su rasgo más distintivo es que no exige saber programar para participar. En lugar de evaluar dominio de lenguajes o herramientas informáticas, propone ejercicios diseñados para que los estudiantes razonen sobre conceptos fundamentales de las ciencias de la computación a partir de situaciones cercanas, visuales o lúdicas. Gracias a este enfoque, Bebras se ha consolidado como una de las experiencias educativas más difundidas para introducir la informática en niveles escolares, mediante una red internacional en la que cada país organiza su propia edición nacional del desafío.

La relevancia de Bebras no se limita a su dimensión pedagógica. Su implementación práctica requiere una infraestructura tecnológica capaz de administrar procesos que van más allá de la simple publicación de ejercicios. Una edición nacional del desafío necesita gestionar participantes, categorías, tareas, pruebas y resultados de manera confiable. En consecuencia, la realización sostenida del concurso depende no solo de criterios educativos, sino también de una plataforma tecnológica que permita operarlo de manera ordenada, segura y escalable.

En distintos países, estas necesidades han dado lugar al desarrollo de plataformas propias, adaptadas a sus contextos organizativos, pedagógicos y técnicos. No existe un único sistema oficial compartido por todos los miembros de la red Bebras, sino una diversidad de implementaciones nacionales que responden a distintas prioridades: algunas se centran en la administración del concurso, otras en la práctica libre de tareas y otras en la comunicación institucional con docentes, familias y coordinadores. Esta diversidad responde al propio modelo de organización de Bebras, que opera mediante estructuras nacionales responsables de ejecutar el desafío en cada país. Por ello, el problema no consiste únicamente en disponer de tareas o materiales, sino en contar con una solución integral capaz de sostener la operación local del desafío.

En Bolivia, esta situación adquiere particular relevancia. A pesar del valor formativo del pensamiento computacional y del crecimiento internacional de Bebras, no se cuenta con una plataforma web propia que permita organizar de forma estable una edición nacional del desafío. Esta carencia resulta especialmente significativa en el contexto de la participación prevista de Bolivia en 2026, ya que la ausencia de infraestructura dificulta la inscripción de participantes, la administración de evaluaciones en línea, la gestión de contenido académico y la consolidación institucional de la iniciativa. Como resultado, la incorporación boliviana a este tipo de experiencias se ve condicionada por barreras operativas que no dependen del interés pedagógico del modelo, sino de la falta de un sistema adecuado para implementarlo en el contexto nacional.

A partir de este problema, el presente trabajo se orienta al desarrollo de la plataforma web del Desafío Bebras Bolivia. La propuesta no se restringe a la construcción de un sitio informativo, sino que plantea un sistema capaz de cubrir los procesos esenciales para la realización del concurso: difusión institucional, inscripción de usuarios, administración de exámenes, resolución de tareas y gestión de ejercicios. De este modo, el proyecto se sitúa en la intersección entre la informática educativa y el desarrollo de software.

Para una persona ajena al tema, es importante entender que este trabajo no busca simplemente digitalizar un examen, sino crear las condiciones técnicas para que una iniciativa educativa internacional pueda existir de forma sostenible en Bolivia. Eso implica comprender, por un lado, qué es el pensamiento computacional y por qué Bebras representa un modelo pertinente para promoverlo y, por otro, qué requerimientos debe cumplir una plataforma web para administrar un concurso educativo de este tipo. Bajo esa lógica, el documento avanza desde los fundamentos conceptuales del desafío hacia el análisis de plataformas de referencia y la definición de una solución tecnológica ajustada al contexto boliviano.

En consecuencia, el presente documento se organiza del siguiente modo. En el capítulo siguiente se desarrollan el marco conceptual y el estado del arte, abordando el pensamiento computacional, la organización internacional Bebras, los criterios para el diseño de tareas y los antecedentes relevantes a nivel internacional y regional. Posteriormente, se expone la metodología adoptada y el marco tecnológico que sustenta la propuesta, incluyendo la retroingeniería de plataformas de referencia y la justificación de las herramientas seleccionadas para la implementación del sistema. Con ello, el trabajo busca fundamentar de manera integral la necesidad y viabilidad de una plataforma web propia para el Desafío Bebras Bolivia.
